\section{`Backyard' models \label{sec:backyard}}

%To recap, the space cost variables to be examined are: distance, time, factors and utility. The first model iteration -  the 'backyard' models - looks at the first three of these. In this model, time is the only factor, and each actor must decide how best to spend their time, splitting it between movement and production. 

\input{ChapterModels/Tables/modelatom_backyard_table}

The `backyard' model is the counterpart to the single-actor love-of-variety trade model outlined in section \ref{sec:modelatom_increts}. That model found that a single actor type, with linear production, could trade over space purely through love of variety. The backyard model is the counterpoint: in a single-actor world, with increasing returns in \emph{production}, what are the minimum conditions that will get them out of their backyard?

The basic choice facing actors is the same: are there more productive options than producing in my own backyard, where output is linear with time input? Every actor is set up such that, if more time is input into their production, increasing returns can be gained. Thus there is an incentive - for example - for two actors to collaborate at one site and split the productive proceeds. 

An immediately obvious difference to the love-of-variety trade model is the nature of the indivisibilities. Whereas trade does not have to happen where goods are produced,actors combining labour must do so at a particular site of production. The barrier in the way of collaboration, then, is the cost for actors to get to the site of production.




There are two elements to the model it is worth discussing in more detail before getting on to the results. Firstly, as mentioned in the table, time is the only factor.
\cmt{model component tables: don't currently list types of good or productive factor. They probably should.}
This is similar to the core model, where labour is the only input. The difference here is that time must be spent between production and space.

Secondly, the model begins with actors of uniform type: there is no initial distinction made between consumers and producers. This is a departure from the standard \gls{ge} assumption of a world where consumers and firms are defined separately. 



